\section{Conclusions}
\label{sec:conclusions}

\subsection{Summary of Contributions}

This thesis investigated how retrieval-augmented large language models can support adaptive maintenance of ETL pipelines under schema evolution. Starting from the observation that even moderate structural changes in source schemas can silently break downstream transformations, the work proposed an integrated workflow that (i) detects and classifies schema evolution events, (ii) represents schemas and changes in machine-readable JSON artefacts, and (iii) leverages an LLM-based mapping regenerator to produce updated transformation code grounded in these artefacts and in existing business logic. The implementation is realised in a Python codebase that operates on both synthetic and real-world inspired datasets, with particular emphasis on a controlled evolution of the NYC yellow taxi schema.

Conceptually, the thesis contributes a concrete instantiation of evidence-grounded retrieval-augmented generation for schema-centric ETL tasks. Rather than treating the LLM as an opaque code generator, the approach constructs a retrieval corpus from historical schema versions, machine-readable change artefacts, and previously validated transformation code, and composes these artefacts with the model to encourage semantic preservation. Practically, the work delivers a reusable schema change detection module, an AI mapping regenerator, and executable ETL pipelines that together form a testbed for studying AI-assisted adaptation workflows. The overall design and implementation bridge ideas from classical schema evolution and mapping adaptation~\cite{velegrakis2003mapping,yu2005semantic,rahm2001survey} with recent advances in LLM-based data engineering~\cite{chen2021evaluating,lewis2020retrieval,narayan2022foundation,buss2025towards}.

\subsection{Answers to the Research Questions}

The empirical results presented in Chapter~5 and interpreted in Chapter~6 allow us to answer the research questions at a high level. For \textbf{RQ1}, the schema change detector is able to recover the designed evolution events exactly in the canonical--V2 scenario, providing a reliable front-end for downstream regeneration under the specific flat tabular schemas considered, while highlighting that some change types (such as removals in other scenarios) remain more challenging. For \textbf{RQ2}, the mapping regenerator can, in favourable cases, produce canonical ingestion mappings that match a strong deterministic baseline exactly (as demonstrated for \texttt{llama3} on NYC taxi V2), while other models exhibit concrete failure modes, underscoring the importance of model choice, prompt design, and validation. For \textbf{RQ3}, the controlled experiments show that the AI-assisted workflow reduces time-to-adaptation relative to a purely manual workflow for the evaluated taxi scenario, without degrading basic data-quality checks, although human review remains essential for production use.

\subsection{Outlook and Future Work}

The thesis opens several directions for future research and development. A first avenue is to scale the evaluation to larger and more heterogeneous corpora of schema histories, including open data portals and versioned enterprise warehouses. This would enable more robust measurements of detection accuracy and adaptation effectiveness, and would allow benchmarking against alternative mapping adaptation tools and industrial ETL platforms.

A second direction is to refine the interaction between schema detection, retrieval, and generation. For example, richer change taxonomies and more sophisticated rename detection could capture complex structural refactorings, while retrieval strategies could be extended to include historical pipeline runs, log messages, or domain documentation. On the generation side, incorporating automated test case synthesis or property-based checks could further reduce the risk of subtle semantic regressions.

Finally, integrating the proposed workflow into production-grade orchestration systems such as Apache Airflow or cloud-native data platforms would make it possible to study continuous adaptation under real operational constraints. This includes exploring human--AI collaboration patterns, designing approval and rollback mechanisms for generated transformations, and investigating how feedback from pipeline monitoring can be fed back into the LLM prompting and retrieval components. Taken together, these directions point towards a future in which ETL pipelines become more resilient and self-adapting, with generative AI acting as a transparent, controllable partner rather than a black-box replacement for human data engineers.
